\documentclass[11pt,letterpaper]{article}

% --- Paquetes de tipografía, idioma y márgenes ---
\usepackage[utf8]{inputenc}
\usepackage[T1]{fontenc}
\usepackage{lmodern}
\usepackage[letterpaper,top=1.8cm,bottom=1.8cm,left=1.8cm,right=1.8cm]{geometry}

% --- Paquetes matemáticos y tablas ---
\usepackage{amsmath,amssymb}
\usepackage{booktabs}
\usepackage{tabularx}
\usepackage{array}
\usepackage{float}
\usepackage{microtype}
\usepackage{xcolor}
\usepackage{graphicx}
\usepackage{fancyhdr}
\usepackage{tcolorbox}
\tcbuselibrary{skins,breakable}
\usepackage{hyperref}
\usepackage{titlesec}
\usepackage{enumitem}
\usepackage{lastpage}
\usepackage{caption}

\setlist{nosep,leftmargin=1.5em}
\setlength{\parindent}{0pt}
\setlength{\parskip}{3.5pt}
\setlength{\abovedisplayskip}{4pt}
\setlength{\belowdisplayskip}{4pt}
\setlength{\abovedisplayshortskip}{2pt}
\setlength{\belowdisplayshortskip}{2pt}
\setlength{\jot}{1.5pt}

% --- Colores Corporativos Solvanta ---
\definecolor{navyblue}{RGB}{26, 79, 124}      % #1a4f7c
\definecolor{forestgreen}{RGB}{20, 120, 40}   % Aprobado
\definecolor{boxbg}{RGB}{250, 251, 253}
\definecolor{calloutblue}{RGB}{240, 245, 250}
\definecolor{borderblue}{RGB}{180, 205, 230}

% --- Formato de Secciones ---
\titleformat{\section}{\large\bfseries\color{navyblue}}{\thesection.}{0.5em}{}
\titleformat{\subsection}{\normalsize\bfseries\color{navyblue}}{\thesubsection.}{0.5em}{}
\titlespacing*{\section}{0pt}{6pt}{2pt}
\titlespacing*{\subsection}{0pt}{4pt}{2pt}

% --- Encabezados y Pies de Página ---
\pagestyle{fancy}
\fancyhf{}
\lhead{\small\color{navyblue}\textbf{Solvanta Engineering Solutions} | DPyPQ}
\rhead{\small Memoria de Cálculo -- Flash HP \texttt{V-501}}
\cfoot{\small Página \thepage\ de \pageref{LastPage}}
\renewcommand{\headrulewidth}{0.4pt}

\hypersetup{
    colorlinks=true,
    linkcolor=navyblue,
    urlcolor=navyblue,
    pdftitle={Memoria de Calculo - Balance de Materia Flash V-501},
    pdfauthor={Solvanta Engineering Solutions}
}

\begin{document}

% --- ENCABEZADO DE TÍTULO ---
\begin{tcolorbox}[colback=navyblue,colframe=navyblue,arc=2mm,coltext=white,boxrule=0pt]
\begin{center}
    {\Large \textbf{MEMORIA TÉCNICA DE CÁLCULO RIGUROSO}}\\[3pt]
    {\large \textbf{Balance de Materia, Masa y Justificación Termodinámica en el Separador Flash HP \texttt{V-501}}}\\[4pt]
    \small \textbf{Documento:} MC-V501-BM01 \quad|\quad \textbf{Unidad:} Bucle de Síntesis de Metanol ($75.0\text{ bar}$, $38.0^\circ\text{C}$) \quad|\quad \textbf{Fecha:} Septiembre 2026
\end{center}
\end{tcolorbox}

\vspace{0.2cm}

% --- RESUMEN EJECUTIVO ---
\begin{tcolorbox}[colback=calloutblue,colframe=borderblue,arc=1mm,title=\textbf{\color{navyblue}Resumen Ejecutivo y Dictamen Físico-Químico},coltitle=navyblue]
\small
\textbf{Pregunta Técnica:} \textit{¿Es físicamente normal que la fracción molar de metanol pase de $0.0192$ ($1.92\%$) en la salida del reactor/enfriador (\texttt{S-404}) a $0.4933$ ($49.33\%$) en el fondo del separador flash (\texttt{S-507})? ¿Se multiplica de la nada o está forzado?}

\vspace{2pt}
\textbf{Dictamen:} \textbf{Es 100\% normal, rigurosamente exacto y termodinámicamente esperado.} 
\begin{itemize}
    \item La fracción molar es una propiedad \textbf{intensiva} que depende del flujo total de la corriente. No se conserva.
    \item La materia se conserva de forma estricta: los moles absolutos de metanol que entran son \textbf{$137.27\text{ kmol/h}$}, y los que salen son exactamente \textbf{$118.15 + 19.12 = 137.27\text{ kmol/h}$}.
    \item El aumento aparente de concentración ocurre porque el flash retira el \textbf{$96.65\%$} del flujo total en forma de gas incondensable ($\text{H}_2$ y $\text{CO}_2$) por el domo, dejando en el fondo una mezcla binaria condensada equimolar de \textbf{metanol crudo y agua} producida en proporción estequiométrica $1:1$.
\end{itemize}
\end{tcolorbox}

\section{Planteamiento Termodinámico y Mecanismo de Separación}

El recipiente \texttt{V-501} es un tambor separador bifásico vertical de alta presión ubicado en el lazo de síntesis.
\begin{itemize}
    \item \textbf{Alimentación (\texttt{S-404}):} Mezcla bifásica proveniente del enfriador \texttt{E-501} a $38.0^\circ\text{C}$ y $75.0\text{ bar}$, con un flujo molar masivo de $7\,149.10\text{ kmol/h}$, dominado por el reciclo de $\text{H}_2$ ($96.03\%$).
    \item \textbf{Domo de Vapor (\texttt{S-501}):} El hidrógeno y el $\text{CO}_2$, debido a su altísima volatilidad relativa y comportamiento supercrítico a estas condiciones, escapan por la parte superior ($6\,909.57\text{ kmol/h}$).
    \item \textbf{Fondos de Líquido (\texttt{S-507}):} El metanol y el agua condensan de forma selectiva y caen al fondo ($239.53\text{ kmol/h}$).
\end{itemize}

\subsection{\texorpdfstring{Demostración Analítica del Factor de Concentración ($25.69\times$)}{Demostración Analítica del Factor de Concentración (25.69x)}}
El factor de reducción de flujo molar entre la entrada y los fondos líquidos es:
\begin{equation}
    \text{Factor de Despojo Gaseoso} = \frac{F}{L} = \frac{7\,149.098\text{ kmol/h}}{239.530\text{ kmol/h}} = 29.846 \approx 29.85
\end{equation}
Considerando una eficiencia de condensación y recuperación del metanol en el líquido del $86.07\%$:
\begin{equation}
    \text{Factor de Concentración} = 29.846 \times 0.86072 = \mathbf{25.689} \approx \mathbf{25.69}
\end{equation}
Multiplicando la fracción inicial por el factor de concentración:
\begin{equation}
    x_{\text{metanol}} = z_{\text{metanol}} \times 25.689 = 0.019201 \times 25.689 = \mathbf{0.4933} \quad (49.33\%)
\end{equation}
Queda demostrado analíticamente que el salto en la fracción es consecuencia matemática de reducir el volumen portador de la fase líquida.

\section{Principio Fundamental de Conservación de la Materia}

La ley universal de conservación de masa para un volumen de control establece:
\begin{equation}
    \text{Entrada} - \text{Salida} + \underbrace{\text{Generación}}_{=\,0} - \underbrace{\text{Consumo}}_{=\,0} = \text{Acumulación}
\end{equation}
Dado que en el tambor flash \texttt{V-501} no ocurren reacciones químicas ($\text{Generación} = 0, \text{Consumo} = 0$) y opera en régimen continuo permanente ($\text{Acumulación} = 0$):
\begin{equation}
    \mathbf{\text{Entrada} - \text{Salida} = \text{Acumulación} = 0 \quad \implies \quad \text{Entrada} = \text{Salida}}
\end{equation}

\clearpage
\section{Cálculo Manual Paso a Paso}

\subsection{\texorpdfstring{Paso 0: Constantes y Pesos Moleculares ($M_i$)}{Paso 0: Constantes y Pesos Moleculares}}
\begin{itemize}[itemsep=1pt]
    \item Hidrógeno ($\text{H}_2$): $M_1 = 2.016\text{ kg/kmol}$ \quad|\quad Dióxido de carbono ($\text{CO}_2$): $M_2 = 44.010\text{ kg/kmol}$
    \item Metanol ($\text{CH}_3\text{OH}$): $M_3 = 32.042\text{ kg/kmol}$ \quad|\quad Agua ($\text{H}_2\text{O}$): $M_4 = 18.015\text{ kg/kmol}$
\end{itemize}

\subsection{\texorpdfstring{Paso 1: Moles de Entrada por Componente ($\dot{n}_{i,\text{in}} = F \times z_i$)}{Paso 1: Moles de Entrada por Componente}}
Para un flujo molar total de entrada $F = 7\,149.098\text{ kmol/h}$:
\begin{align*}
    \dot{n}_{\text{H}_2,\text{in}} &= 7\,149.098 \times 0.96025 = \mathbf{6\,864.93\text{ kmol/h}} &
    \dot{n}_{\text{CO}_2,\text{in}} &= 7\,149.098 \times 0.00228 = \mathbf{16.33\text{ kmol/h}} \\
    \dot{n}_{\text{MeOH},\text{in}} &= 7\,149.098 \times 0.01920 = \mathbf{137.27\text{ kmol/h}} &
    \dot{n}_{\text{H}_2\text{O},\text{in}} &= 7\,149.098 \times 0.01826 = \mathbf{130.57\text{ kmol/h}}
\end{align*}
\textbf{Flujo molar total de entrada:}\quad $\sum \dot{n}_{i,\text{in}} = 6\,864.93 + 16.33 + 137.27 + 130.57 = \mathbf{7\,149.10\text{ kmol/h}}$.

\subsection{\texorpdfstring{Paso 2: Masa de Entrada por Componente ($\dot{m}_{i,\text{in}} = \dot{n}_{i,\text{in}} \times M_i$)}{Paso 2: Masa de Entrada por Componente}}
\begin{align*}
    \dot{m}_{\text{H}_2,\text{in}} &= 6\,864.93 \times 2.016 = \mathbf{13\,839.70\text{ kg/h}} &
    \dot{m}_{\text{CO}_2,\text{in}} &= 16.33 \times 44.010 = \mathbf{718.55\text{ kg/h}} \\
    \dot{m}_{\text{MeOH},\text{in}} &= 137.27 \times 32.042 = \mathbf{4\,398.59\text{ kg/h}} &
    \dot{m}_{\text{H}_2\text{O},\text{in}} &= 130.57 \times 18.015 = \mathbf{2\,352.15\text{ kg/h}}
\end{align*}
\textbf{Flujo másico total de entrada:}\quad $\sum \dot{m}_{i,\text{in}} = \mathbf{21\,308.08\text{ kg/h}}$ (Consistente al $100\%$ con DWSIM).

\subsection{\texorpdfstring{Paso 3: Separación de Fases en \texttt{V-501} ($V$ y $L$)}{Paso 3: Separación de Fases en V-501}}
A $38.0^\circ\text{C}$ y $75.0\text{ bar}$, la relación de fases calculada por equilibrio riguroso es:
\begin{equation*}
    V = 6\,909.568\text{ kmol/h} \quad (\text{Domo \texttt{S-501}}), \qquad L = 239.530\text{ kmol/h} \quad (\text{Fondos \texttt{S-507}})
\end{equation*}
\textbf{A. Distribución Molar en Domo de Vapor (\texttt{S-501}):}
\begin{align*}
    \dot{n}_{\text{H}_2,V} &= 6\,909.57 \times 0.99348 = \mathbf{6\,864.52\text{ kmol/h}} &
    \dot{n}_{\text{CO}_2,V} &= 6\,909.57 \times 0.00235 = \mathbf{16.23\text{ kmol/h}} \\
    \dot{n}_{\text{MeOH},V} &= 6\,909.57 \times 0.00277 = \mathbf{19.12\text{ kmol/h}} &
    \dot{n}_{\text{H}_2\text{O},V} &= 6\,909.57 \times 0.00140 = \mathbf{9.70\text{ kmol/h}}
\end{align*}
\textbf{B. Distribución Molar en Fondos de Líquido (\texttt{S-507}):}
\begin{align*}
    \dot{n}_{\text{H}_2,L} &= 239.53 \times 0.00172 = \mathbf{0.41\text{ kmol/h}} &
    \dot{n}_{\text{CO}_2,L} &= 239.53 \times 0.00040 = \mathbf{0.10\text{ kmol/h}} \\
    \dot{n}_{\text{MeOH},L} &= 239.53 \times 0.49327 = \mathbf{118.15\text{ kmol/h}} &
    \dot{n}_{\text{H}_2\text{O},L} &= 239.53 \times 0.50461 = \mathbf{120.87\text{ kmol/h}}
\end{align*}

\subsection{\texorpdfstring{Paso 4: Comprobación de Acumulación ($\text{Entrada}_i - \text{Salida}_i = \text{Acumulación}_i$)}{Paso 4: Comprobación de Acumulación}}
\begin{align*}
    \text{Acum}(\text{H}_2) &= 6\,864.93 - 6\,864.93 = \mathbf{0.00\text{ kmol/h}} &
    \text{Acum}(\text{CO}_2) &= 16.33 - 16.33 = \mathbf{0.00\text{ kmol/h}} \\
    \text{Acum}(\text{MeOH}) &= 137.27 - 137.27 = \mathbf{0.00\text{ kmol/h}} &
    \text{Acum}(\text{H}_2\text{O}) &= 130.57 - 130.57 = \mathbf{0.00\text{ kmol/h}}
\end{align*}

\subsection{Paso 5: Masa que Sale y Cierre Global de Masa}
\begin{align*}
    \dot{m}_{\text{Vapor}} &= 13\,838.04 + 714.31 + 612.66 + 174.74 = \mathbf{15\,339.75\text{ kg/h}} \\
    \dot{m}_{\text{Líquido}} &= 0.83 + 4.23 + 3\,785.86 + 2\,177.47 = \mathbf{5\,968.39\text{ kg/h}}
\end{align*}
\textbf{Salidas Totales:} $15\,339.75 + 5\,968.39 = \mathbf{21\,308.14\text{ kg/h}}$ \quad|\quad \textbf{Acumulación:} $\mathbf{-0.06\text{ kg/h}}$ ($0.0003\%$).

\clearpage
\section{Tablas Oficiales Consolidadas de Balance}

\subsection{Tabla 1: Balance Molar Completo y Comprobación de Acumulación}
\begin{table}[H]
\centering
\small
\setlength{\tabcolsep}{3pt}
\begin{tabularx}{\textwidth}{>{\raggedright\arraybackslash}p{3.5cm} >{\raggedleft\arraybackslash}X >{\centering\arraybackslash}p{1.6cm} >{\raggedleft\arraybackslash}X >{\raggedleft\arraybackslash}X >{\raggedleft\arraybackslash}X >{\raggedleft\arraybackslash}p{2.3cm}}
\toprule
\textbf{Componente} & \textbf{Entrada S-404} & \textbf{Frac. $z_i$} & \textbf{Domo S-501} & \textbf{Fondos S-507} & \textbf{Salidas} & \textbf{Acumulación} \\
 & ($\text{kmol/h}$) & (Molar) & ($\text{kmol/h}$) & ($\text{kmol/h}$) & ($\text{kmol/h}$) & ($\text{kmol/h}$) \\
\midrule
Hidrógeno ($\text{H}_2$) & $6\,864.932$ & $0.9603$ & $6\,864.517$ & $0.412$ & $6\,864.929$ & $\mathbf{+0.003}$ \\
Dióxido de Carbono ($\text{CO}_2$) & $16.327$ & $0.0023$ & $16.231$ & $0.096$ & $16.327$ & $\mathbf{0.000}$ \\
Metanol ($\text{CH}_3\text{OH}$) & $\mathbf{137.273}$ & $\mathbf{0.0192}$ & $19.121$ & $\mathbf{118.154}$ & $\mathbf{137.275}$ & $\mathbf{-0.002}$ \\
Agua ($\text{H}_2\text{O}$) & $130.566$ & $0.0183$ & $9.700$ & $120.868$ & $130.568$ & $\mathbf{-0.002}$ \\
\midrule
\textbf{TOTAL} & $\mathbf{7\,149.098}$ & $\mathbf{1.0000}$ & $\mathbf{6\,909.568}$ & $\mathbf{239.530}$ & $\mathbf{7\,149.098}$ & $\mathbf{0.000}$ \\
\bottomrule
\end{tabularx}
\caption{Balance de materia molar riguroso en el tambor flash \texttt{V-501}.}
\end{table}

\vspace{-0.1cm}
\subsection{Tabla 2: Balance Másico Completo y Comprobación de Acumulación}
\begin{table}[H]
\centering
\small
\setlength{\tabcolsep}{3pt}
\begin{tabularx}{\textwidth}{>{\raggedright\arraybackslash}p{3.5cm} >{\raggedleft\arraybackslash}X >{\centering\arraybackslash}p{1.6cm} >{\raggedleft\arraybackslash}X >{\raggedleft\arraybackslash}X >{\raggedleft\arraybackslash}X >{\raggedleft\arraybackslash}p{2.3cm}}
\toprule
\textbf{Componente} & \textbf{Entrada S-404} & \textbf{Frac. $w_{\text{in}}$} & \textbf{Domo S-501} & \textbf{Fondos S-507} & \textbf{Salidas} & \textbf{Acumulación} \\
 & ($\text{kg/h}$) & (Masa) & ($\text{kg/h}$) & ($\text{kg/h}$) & ($\text{kg/h}$) & ($\text{kg/h}$) \\
\midrule
Hidrógeno ($\text{H}_2$) & $13\,838.88$ & $64.95\%$ & $13\,838.04$ & $0.83$ & $13\,838.87$ & $\mathbf{+0.01}$ \\
Dióxido de Carbono ($\text{CO}_2$) & $718.54$ & $3.37\%$ & $714.31$ & $4.23$ & $718.54$ & $\mathbf{0.00}$ \\
Metanol ($\text{CH}_3\text{OH}$) & $\mathbf{4\,398.47}$ & $\mathbf{20.64\%}$ & $612.66$ & $\mathbf{3\,785.86}$ & $\mathbf{4\,398.52}$ & $\mathbf{-0.05}$ \\
Agua ($\text{H}_2\text{O}$) & $2\,352.19$ & $11.04\%$ & $174.74$ & $2\,177.47$ & $2\,352.21$ & $\mathbf{-0.02}$ \\
\midrule
\textbf{TOTAL} & $\mathbf{21\,308.08}$ & $\mathbf{100.00\%}$ & $\mathbf{15\,339.75}$ & $\mathbf{5\,968.39}$ & $\mathbf{21\,308.14}$ & $\mathbf{-0.06}$ \\
\bottomrule
\end{tabularx}
\caption{Balance másico riguroso y error de cierre en el tambor flash \texttt{V-501}.}
\end{table}

\vspace{-0.1cm}
\subsection{Tabla 3: Desempeño y Eficiencias de Separación de Fases}
\begin{table}[H]
\centering
\small
\begin{tabularx}{\textwidth}{>{\raggedright\arraybackslash}p{4.6cm} >{\centering\arraybackslash}p{3.2cm} >{\centering\arraybackslash}p{1.4cm} >{\raggedright\arraybackslash}X}
\toprule
\textbf{Parámetro} & \textbf{Valor} & \textbf{Unidad} & \textbf{Interpretación Técnica} \\
\midrule
Recuperación de Metanol en Líquido ($\eta_{\text{MeOH}}$) & $\mathbf{86.07\%}$ & $\%$ & Cosecha líquida de metanol crudo enviada a desgasificador \texttt{V-601}. \\
Arrastre de Metanol en Vapor & $13.93\%$ & $\%$ & Metanol recirculado al reactor vía compresor \texttt{C-501}. \\
Recuperación de Agua en Líquido ($\eta_{\text{Agua}}$) & $92.57\%$ & $\%$ & Retención de agua en fondos para posterior destilación. \\
Separación de Hidrógeno en Domo ($\eta_{\text{H}_2}$) & $\mathbf{99.99\%}$ & $\%$ & Purificación del gas portador para reciclo sin carga inútil. \\
Composición Líquida Metanol / Agua & $\mathbf{49.33\%}$ / $\mathbf{50.46\%}$ & $\%$ mol & Relación equimolar estequiométrica ($\text{CO}_2 + 3\text{H}_2 \to \text{CH}_3\text{OH} + \text{H}_2\text{O}$). \\
\bottomrule
\end{tabularx}
\caption{Indicadores clave de eficiencia termodinámica de \texttt{V-501}.}
\end{table}

\clearpage
\section{Dictamen y Conclusiones Finales}

\begin{enumerate}[itemsep=4pt]
    \item \textbf{Sin Creación de Materia:} En ningún momento se multiplica el metanol. Los $137.27\text{ kmol/h}$ que entran al tambor flash se conservan exactamente entre el domo y el fondo ($19.12 + 118.15 = 137.27\text{ kmol/h}$).
    \item \textbf{Causa del Incremento de Concentración ($1.92\% \to 49.33\%$):} El aumento en la fracción molar se debe a la eliminación del $96.65\%$ de los moles totales como vapor incondensable en el domo. Al reducirse el divisor total de la mezcla líquida a solo $239.53\text{ kmol/h}$, la proporción de metanol aumenta por un factor de $25.69$.
    \item \textbf{Consistencia de la Relación Equimolar:} El hecho de que el fondo contenga $\approx 49.33\%$ de metanol y $\approx 50.46\%$ de agua valida que la reacción de síntesis está operando estequiométricamente según la reacción de hidrogenación directa de $\text{CO}_2$, condensando ambos productos juntos.
    \item \textbf{Validez Numérica:} La acumulación calculada tanto molar ($0.000\text{ kmol/h}$) como másica ($-0.06\text{ kg/h}$ frente a $21\,308.08\text{ kg/h}$) demuestra un cierre perfecto del $99.9997\%$, cumpliendo con los más estrictos estándares de ingeniería de procesos.
\end{enumerate}

\vspace{0.4cm}
\begin{center}
    \small\color{navyblue}\textbf{Solvanta Engineering Solutions} -- Departamento de Procesos y Simulación Numérica\\
    \textit{Documento generado y verificado contra modelos termodinámicos rigurosos de DWSIM v8.8.}
\end{center}

\end{document}
